\hypertarget{appendix-505a502c372a}{}\label{referencestart:appendix-505a502c372a}
\section{Projection rejection and scalar components}\label{appendix-505a502c372a-projection-rejection-and-scalar-components}

\textbf{Curriculum depth:} Introduced · \textbf{First seen:} Phase6 · \textbf{Latest development:} Phase6

\subsection{Definition and mechanism}\label{appendix-505a502c372a-definition-and-mechanism}

Projection onto a nonzero axis is the along-axis vector whose remainder is orthogonal. Rejection is the remainder. Signed scalar projection measures length; the projection coefficient multiplies the target axis and need not equal that length.

For input \((3,4,0)\) and target \((2,0,0)\), coefficient is \(1.5\), scalar projection \(3\), projection \((3,0,0)\) and rejection \((0,4,0)\).

\subsection{Prerequisites}\label{appendix-505a502c372a-prerequisites}

\begin{itemize}
\tightlist
\item
  \hyperref[appendix-8bc4e8fc547c]{Dot products and orthogonality} --- Inner-product measurements supply the perpendicularity and alignment constraints.
\end{itemize}

\subsection{Used by}\label{appendix-505a502c372a-used-by}

\begin{itemize}
\item
  \hyperref[appendix-e31893dec985]{Surface tangents normals and bitangents} --- Builds on this mechanism in the connected geometry story.
\item
  \hyperref[appendix-44e00863b4d5]{Reflection refraction and shading directions} --- Builds on this mechanism in the connected geometry story.
\item
  \hyperref[appendix-8572a4a060c1]{Linear systems and fundamental subspaces} --- Extends this foundation into linear-systems and decomposition analysis.
\item
  \hyperref[appendix-3c73f7862157]{Exceptional projections and dimension} \#\# Recall and next depth
\end{itemize}

Derive the denominator, closest-point property and idempotence without assuming a unit target. Full production proofs, target-specific integration and learner mastery require further evidence.

\subsection{Source and navigation}\label{appendix-505a502c372a-source-and-navigation}

\href{../02-Mathematics-for-Visual-Computing/main.pdf\#nameddest=01-foundations-i-ch017-projection-separates-useful-movement-from-sideways-movement}{Teaching narrative (Mathematics for Visual Computing, Mathematics for Visual Computing, Chapter 1)} · \hyperref[appendix-fd090d982a15-phase-6---parts-501-600]{Coverage and source gap} · \hyperref[appendix-a103e7846d21-phase-6---source-analysis-and-verification]{Evidence and limits} · \hyperref[appendix-4ad47bdd269a]{Dependency Map}
